\section{Analytic Geometry, Condensed Mathematics, and Analytic Rings}
\label{lec:1}

This first lecture is an introduction to the first half of the course, whose
goal is to set up foundations for analytic geometry: a general theory of
\emph{analytic stacks}. Clausen began with motivation, drawn from the existing
classical theories of analytic geometry and their shortcomings; then introduced
the condensed formalism as a replacement for topological spaces; and finally
gave the notion of \emph{analytic ring}, the affine local models of the theory,
closing with the example of solid analytic rings. Throughout, this is an
introduction: the definitions are stated in outline, and the details are
deferred to the coming lectures.

\subsection{Motivation: classical theories of analytic geometry}

Classically there are several different theories of analytic geometry. Clausen
listed the ones he is familiar with, with the caveat that the list is not
exhaustive.

\begin{enumerate}
\item \emph{Complex analytic spaces.} In the smooth case these are the complex
manifolds: spaces obtained by gluing open subsets of $\mathbb{C}^n$ along
biholomorphisms, i.e.\ maps that locally admit a power series expansion. In the
non-smooth case one also allows, locally, the zero locus of a finite collection
of analytic functions on such open subsets. This is perhaps the gold standard.

\item \emph{Locally analytic manifolds.} This is a generalization of the smooth
case, presented in Serre's book on Lie groups and Lie algebras
\cite{Serre_1992}. One starts with a complete normed field $(k,\lvert\cdot\rvert)$,
which may be Archimedean or non-Archimedean, and glues open subsets of $k^d$
along locally analytic isomorphisms -- maps admitting, locally around every
point, a convergent power series expansion with coefficients in $k$. When
$k=\mathbb{C}$ this recovers the smooth case of (1); when $k=\mathbb{R}$ it is a
version of the theory of real manifolds; both are reasonable theories. But when
$k$ is non-Archimedean the theory is not geometrically rich, absent some extra
structure such as a group structure.
\end{enumerate}

\begin{remark}[Total disconnectedness in the non-Archimedean case]
\label{rmk:1:totally-disconnected}
The failure of richness for, say, $k=\mathbb{Q}_p$ comes from the topology of
$\mathbb{Q}_p^d$ being totally disconnected: every unit ball breaks up into $p$
smaller unit balls, each of which breaks up again, and so on. Serre's book
\cite{Serre_1992} contains the classification of compact $p$-adic manifolds; it
is a simple combinatorial classification (by dimension and one further
invariant), reflecting how little geometry is present.
\end{remark}

The next two theories react to this by shifting attention from local topology to
local rings of functions.

\begin{enumerate}
\setcounter{enumi}{2}
\item \emph{Rigid analytic geometry}, introduced by Tate \cite{Tate_1971}, over
a non-Archimedean field $k$ -- motivated in part by the uniformization of
elliptic curves. This is a geometrically rich theory that works in the
non-Archimedean case. In contrast to (1)--(2) one does not think in terms of
gluing open subsets of $k^d$, or even in topological terms at all, but rather
algebraically. The shift is twofold:
\begin{itemize}
\item one focuses on the \emph{local rings of functions} rather than on local
topology, letting the ring of functions dictate what the topology should be; and
\item the local models are rings of functions convergent not on an \emph{open}
disk but on a \emph{closed} disk -- which is what makes sense in the
non-Archimedean setting.
\end{itemize}
Concretely the local rings of functions are (quotients of) rings of functions
convergent on a closed polydisk; these are the Tate algebras. The manner of
gluing these local models to global rigid analytic spaces is halfway between
algebraic geometry and ordinary analytic geometry.
\end{enumerate}

\begin{remark}[Terminology]
\label{rmk:1:affinoid-terminology}
There was a brief exchange about whether ``Tate algebra'' should mean the rings
of functions convergent on a polydisk, or the quotients of these; the usage
varies between references. Clausen settled on calling the quotients
\emph{affinoid algebras}.
\end{remark}

\begin{enumerate}
\setcounter{enumi}{3}
\item \emph{Adic spaces}, Huber's generalization of rigid analytic geometry
\cite{Huber_1994}. Rigid analytic geometry takes place over a fixed base field;
a loose analogy is that adic spaces are to rigid analytic spaces as general
schemes are to varieties over a field. One again specifies local models
algebraically, by a ring of functions, but now with a new twist: the datum is a
pair
\[
A^{+}\subseteq A,
\]
a topological ring $A$ (``local ring of functions'', not local in the sense of
commutative algebra) together with a subring $A^{+}$. To $A$ one attaches a
space of valuations, and $A^{+}$ singles out those valuations for which the
elements of $A^{+}$ are regarded as integral. A useful flexibility of Huber's
theory is that one may consider different choices of $A^{+}$ for a fixed $A$. We
will see later in the lecture what this choice is really doing.

\item \emph{Berkovich's theory.} The rings of functions in the rigid and adic
settings are typically rings complete with respect to a finitely generated
ideal, given the inverse-limit topology (all quotients discrete), possibly with
something inverted -- provided the inversion inverts what one completed along.
For example, functions on the closed unit disk in $\mathbb{A}^1$ over
$\mathbb{Q}_p$ arise as
\[
\bigl(\mathbb{Z}[T]^{\wedge}_{p}\bigr)\!\left[\tfrac{1}{p}\right],
\]
the $p$-adic completion of $\mathbb{Z}[T]$ with $p$ inverted. These are $p$-adic
Banach rings. Berkovich instead works with arbitrary Banach rings
$(R,\lvert\cdot\rvert)$ as local models, and attaches to such a ring the space
$\mathcal{M}(R,\lvert\cdot\rvert)$ of multiplicative seminorms, a compact
Hausdorff space; the allowed gluings are organized by this Berkovich space. Since
$\mathbb{R}$ and $\mathbb{C}$ are Banach rings, this theory -- alone among those
listed -- allows both Archimedean and non-Archimedean phenomena. Complex
analytic spaces, and in a certain sense rigid analytic geometry, appear as
special cases. The global theory, however, does not function as smoothly as in
the other settings; see \cite{Berkovich_2012}.
\end{enumerate}

The relationships among these theories are more or less well understood: (1) is a
special case of (2) and (of the smooth part of) others, rigid geometry is a
special case of adic spaces and of Berkovich's theory, and so on. One can
formulate the comparisons by hand, but there is no common framework in which all
of these examples live and in which the comparisons take place.

\subsection{Why introduce a new theory?}

Clausen gave several reasons, theoretical and practical.

\begin{itemize}
\item \emph{To accommodate all examples.} One would like a single general theory
of analytic spaces that can be specialized to whichever context one is interested
in, rather than a list of separate theories with their own flavors.

\item \emph{Archimedean and non-Archimedean geometry together, with good
gluing.} Of the theories above, the only rich one allowing both Archimedean and
non-Archimedean geometry is Berkovich's, but there the gluing is not so well
worked out. Berkovich restricted essentially to the non-Archimedean case from the
start; more general gluings were investigated by Poineau. In these approaches one
fixes a Banach ring as a base, defines affine space of some dimension over it,
and glues along certain subsets of affine space. The objects one glues along are
often not themselves controlled by Banach rings, so the nature of the gluing is
constrained to a finite-type situation and does not fit the mould of how one
passes from affine schemes to general schemes by naive gluing of local models.

\item \emph{Flexibility, and issues of descent.} Even within their own contexts
these theories are less flexible than the theory of schemes, largely for reasons
of descent. For a scheme one has the category of quasi-coherent sheaves
$\QCoh(-)$: over a ring one takes all modules, and this glues. In analytic
geometry there is no such theory in any of the classical settings. One always has
a local ring $R$ describing the local geometry, and one may assign to it the
category of all $R$-modules, but this does not glue: modules on two pieces with
gluing data on the overlap need not globalize. Only finitely presented modules
glue, giving the theory of coherent sheaves $\Coh(-)$ -- a beautiful and
central tool, but constrained by finiteness. For instance, the pushforward of the
structure sheaf along a map of analytic spaces is not coherent unless the map is
finite, whereas such pushforwards are basic examples of quasi-coherent sheaves in
algebraic geometry, where no finiteness is built in.

\item \emph{A speculative, practical reason from the Langlands program.} Fargues
and Scholze geometrized the local Langlands correspondence
\cite{fargues2024geometrizationlocallanglandscorrespondence}, clarifying the
local Langlands program. The geometrization is based on replacing $\mathbb{Q}_p$
by a more exotic object, the Fargues--Fontaine curve -- or rather a family of
such curves -- produced in the language of adic spaces: an adic space over
$\mathbb{Q}_p$, not at all of finite type, which the theory of adic spaces
fortunately existed to accommodate. Speculatively, one might hope to geometrize
not just local but global Langlands. (Clausen noted that Scholze is optimistic
about this, and he himself less so.) This would involve replacing $\mathbb{Q}$
or $\mathbb{Z}$ by some family of exotic analytic spaces -- which would have to
have both Archimedean and non-Archimedean aspects (there should also be a version
over $\mathbb{R}$, which Clausen believes Scholze is working out) and would not
be of finite type in any sense. No existing theory could give the language to
describe such an object, if it exists; and it is valuable to have a precise
theory to guide the exploration of such possibilities.
\end{itemize}

The goal of the course is thus to introduce a new theory of analytic geometry and
to explain its relation to the previous theories. Today's lecture treats only the
affine situation -- the analytic rings.

\subsection{Condensed mathematics}

The descent problems above can be traced to the fact that the local rings in all
these theories are not abstract rings but \emph{topological} rings, and it is
essential to remember the topology.

\begin{example}[Completed vs.\ algebraic tensor products]
\label{ex:1:completed-tensor}
In algebraic geometry the polynomial ring in two variables, functions on
$\mathbb{A}^2$, is the tensor product $\mathbb{C}[x]\otimes_{\mathbb{C}}\mathbb{C}[y]$.
In rigid analytic geometry, however, the algebraic tensor product of the
one-variable Tate algebra with itself gives a wrong, ``crazy'' ring, because it
forgets the $p$-adic topology; one must take the $p$-adically completed tensor
product to obtain the correct ring of functions in two variables. Tensor products
compute fiber products geometrically, so one must remember the topology when
forming them.
\end{example}

This is the source both of the absence of a naive theory of quasi-coherent
sheaves and of the gluing problems outside the finitely generated case. Yet
topological modules over a topological ring are not suitable for a general
theory: the category of topological modules over a topological ring is not
abelian, however one dresses it up.

\begin{remark}[Why topological modules fail to be abelian]
\label{rmk:1:not-abelian}
The phenomenon is that a dense inclusion of modules -- which occurs constantly
in infinite-dimensional situations -- is generally both an epimorphism and a
monomorphism in reasonable such categories, without being an isomorphism. (One
must impose a separatedness/Hausdorff condition for the literal statement; the
underlying point is that a non-strict map has an image carrying two different
topologies.)
\end{remark}

The remedy is to go back to basics and define a replacement for the category of
topological spaces, one on which it is easy to pile algebraic structure --
analogues of topological rings and modules -- so as to obtain an abelian
category in the end. The basic idea is old (Clausen attributed frequent use of it
to Grothendieck, and suspected it is older still).

The difficulty with a topological space is that it is second-order data: a set of
points together with a set of subsets of it, which is what makes topology hard to
mix with algebraic structure. One would rather record first-order data. The idea
is to single out a collection of nice \emph{test spaces} $S$, and to encode a
space $X$ not as a set with a chosen family of open subsets, but by recording the
data of the (would-be) continuous maps $S\to X$, axiomatizing the structure and
properties visible in that situation. We care about $X$ only insofar as it is
seen by maps from the test objects.

\begin{remark}[Older incarnations]
\label{rmk:1:older-incarnations}
Traditionally, to do homotopy theory, one sometimes fixed a set $X$ together with
a collection of maps of simplices into $X$ and defined homotopy and homology from
that. Clausen tentatively named Spanier as a source; he was not certain of the
attribution.
\end{remark}

We now fix the test spaces to be profinite sets, and impose a countability
restriction.

\begin{definition}[Light profinite set]
\label{def:1:light-profinite}
A \emph{light profinite set} is a countable inverse limit of finite sets (each
finite set discrete, the limit given the inverse-limit topology). Equivalently,
it is a metrizable totally disconnected compact Hausdorff space; equivalently
still, a compact Hausdorff space with a countable basis of clopen subsets. We
write $\LightProf$ for the category of light profinite sets, morphisms being
continuous maps (equivalently, maps of the associated countable pro-systems of
finite sets).
\end{definition}

\begin{remark}[Why ``light'']
\label{rmk:1:why-light}
The first courses on condensed mathematics (Scholze's, following
\cite{clausen2019condensed}) used arbitrary profinite sets, i.e.\ arbitrary
totally disconnected compact Hausdorff spaces, as test spaces. That is a large
category with no cardinal bound, so encoding the data below records more than a
set's worth of information, causing technical trouble which was worked around.
For this course we take instead the light variant; the reasons for this precise
choice will be explained over the first few lectures.
\end{remark}

\begin{definition}[Light condensed set]
\label{def:1:light-condensed}
A \emph{light condensed set} is a sheaf of sets on $\LightProf$ for the
Grothendieck topology in which a covering sieve is one containing a finite,
jointly surjective collection of continuous maps. Thus finite disjoint unions are
covers, and a surjection is a cover.
\end{definition}

Unwinding the sheaf condition makes this elementary.

\begin{definition}[Light condensed set, explicitly]
\label{def:1:light-condensed-explicit}
A light condensed set is a functor $X\colon (\LightProf)^{\mathrm{op}}\to\mathrm{Set}$
such that:
\begin{enumerate}
\item $X(\emptyset)=\ast$;
\item $X(S\amalg T)\xrightarrow{\ \sim\ } X(S)\times X(T)$; and
\item for every surjection $T\twoheadrightarrow S$,
\begin{equation*}
X(S)\xrightarrow{\ \sim\ }\eq\bigl(X(T)\rightrightarrows X(T\times_S T)\bigr).
\end{equation*}
\end{enumerate}
\end{definition}

\begin{example}[Topological spaces as condensed sets]
\label{ex:1:top-to-condensed}
Any topological space $X$ gives a light condensed set with functor of points
\begin{equation*}
S\longmapsto X(S):=\Cont(S,X),
\end{equation*}
the set of continuous maps $S\to X$. Functoriality and the three properties of
\Cref{def:1:light-condensed-explicit} are immediate.
\end{example}

\begin{remark}[Grothendieck topologies will be used in earnest]
\label{rmk:1:use-grothendieck}
Although the definition can be unwound as above, the elementary appearance is
somewhat illusory: the theory of Grothendieck topologies and sheaves is used
seriously in this course, both here and later in gluing the affine (analytic
ring) case to general analytic spaces. Clausen recommended that anyone unfamiliar
with sheaves on sites read up on the subject before following the course.
\end{remark}

Two remarks indicate why this precise definition is made.

\begin{remark}[The two basic test objects]
\label{rmk:1:two-test-objects}
The two most important light profinite sets are the point $\ast$ and the
one-point compactification $\mathbb{N}\cup\{\infty\}$ of the natural numbers.
Evaluating at $\ast$ recovers the underlying set of a condensed set, and
evaluating at $\mathbb{N}\cup\{\infty\}$ records its convergent sequences (for a
topological space these are literally the underlying set and the set of
convergent sequences; in general they are the abstract analogues).
\end{remark}

\begin{remark}[Effect of the two choices in the topology]
\label{rmk:1:covers-and-finitary}
Two features of the Grothendieck topology pull in complementary directions.
\begin{itemize}
\item Allowing \emph{all} surjections to count as covers simplifies the structure
of the category, and in particular yields good homological algebra when one
passes to light condensed abelian groups: one has much flexibility to work
locally.
\item Requiring the topology to be \emph{finitary} (covers generated by finite
jointly surjective families) gives good categorical compactness properties for
the light profinite sets, sitting by the Yoneda embedding inside all light
condensed sets -- and indeed for all metrizable compact Hausdorff spaces. For
example, the unit interval has the surjection
\begin{equation*}
\{0,1\}^{\mathbb{N}}\twoheadrightarrow [0,1]
\end{equation*}
from the Cantor set given by binary expansions; since $\{0,1\}^{\mathbb{N}}$ is
one of the basic test objects and the topology is finitary, the topological
compactness of these compact Hausdorff spaces translates into a categorical
compactness property inside the larger category. For this it is important to have
the larger light profinite sets available, not merely the sets of convergent
sequences.
\end{itemize}
\end{remark}

\begin{notation}[Dropping ``light'']
\label{not:1:drop-light}
From now on \emph{condensed set} means light condensed set and \emph{profinite
set} means light profinite set; the countability hypothesis is always present.
\end{notation}

\subsection{Analytic rings}

Condensed sets give rise, immediately from the sheaf-theoretic viewpoint, to
condensed rings and condensed modules: a \emph{condensed ring} is a sheaf of
(commutative) rings on the site of profinite sets, equivalently a commutative
ring object in condensed sets, and a condensed module over it is a module object.

\begin{remark}
\label{rmk:1:condensed-ring-pointwise}
Concretely a condensed ring is a compatible family of abstract commutative rings
$R(S)$, one for each profinite set $S$, satisfying the conditions of
\Cref{def:1:light-condensed-explicit} (so $R(S\amalg T)=R(S)\times R(T)$ is a
\emph{Cartesian} product of rings, etc.). ``Ring'' means commutative ring
throughout.
\end{remark}

One might hope to use condensed rings directly as the local models of analytic
geometry, by analogy with schemes (which are built from discrete rings). This is
not enough: condensed rings alone do not give a good geometry.

\begin{remark}[Why condensed rings alone are insufficient]
\label{rmk:1:condensed-rings-insufficient}
The category of condensed rings has pushouts given by relative tensor products,
just as for classical commutative rings, and these should geometrically compute
fiber products -- i.e.\ should be the completed tensor products of
\Cref{ex:1:completed-tensor}. But for condensed rings $A,B$ over $K$ one computes,
from the nature of the Grothendieck topology,
\begin{equation*}
\bigl(A\otimes_{K} B\bigr)(\ast)=A(\ast)\otimes_{K(\ast)} B(\ast),
\end{equation*}
so the underlying ring of the relative tensor product is the \emph{abstract}
(uncompleted) tensor product of underlying rings. The relative tensor product
thus merely puts a (nontrivial) condensed structure on the abstract tensor
product; it does not perform any completion, and completion does change the
underlying set. So this is not yet doing the geometrically correct thing.
\end{remark}

To fix this, one puts additional structure on a condensed ring, recording a class
of condensed modules to be considered as \emph{complete} in an appropriate sense.
This is the notion of analytic ring: a condensed ring together with data telling
which condensed modules are complete for the theory being described.

Before the definition, one must be precise about ``ring''.

\begin{remark}[Which notion of ring: animated]
\label{rmk:1:which-ring}
Experience in algebraic geometry shows that the correct fiber product of schemes
is the \emph{derived} fiber product, corresponding on affines to the derived
relative tensor product of rings. Thanks to Lurie's work \cite{lurie2017higher}
this is no longer a technical hassle, so the course works with derived rings. In
practice almost all basic examples will carry no derived structure -- they will
be ordinary rings -- so one can follow the course while imagining everything is
an ordinary ring; but the general claims can fail if everything is taken to be
ordinary.

Once one works derived there are inequivalent choices. The two basic options are
$\mathbb{E}_\infty$-algebras (in the sense of spectra) and \emph{animated}
commutative rings, i.e.\ the objects presented by simplicial commutative rings;
and in each case one may or may not allow negative homotopy. We do not allow
negative homotopy (rings are connective), and we choose animated commutative
rings, which are more directly tied to classical algebraic geometry: taking
derived tensor products of ordinary rings always lands in this class. The whole
theory works equally well in any of the variants -- and is in fact slightly less
technical to set up in this seemingly more complicated setting.
\end{remark}

\begin{definition}[Condensed animated ring]
\label{def:1:condensed-animated-ring}
A \emph{condensed animated ring} is a hypersheaf of animated rings on the site of
(light) profinite sets. From now on ``ring'' means animated ring, and if we wish
to stress that a ring lives in degree zero we call it \emph{static} (or
classical).
\end{definition}

\begin{remark}[Why hypersheaves]
\label{rmk:1:hypersheaves}
Hypersheaves are used (rather than sheaves) because one wants, e.g., convergence
of Postnikov towers, which can be proved for hypersheaves but is not known for
sheaves; one cannot show sheaves and hypersheaves agree. Moreover everything can
always be shown to be a hypersheaf, and in practice hypercovers cause no more
trouble than ordinary covers, so nothing is lost.
\end{remark}

\begin{definition}[Derived category of a condensed animated ring]
\label{def:1:derived-category}
The basic invariant of a condensed animated ring $R$ is its full derived category
$D(R)$: the (unbounded) $\infty$-category of hypersheaves of modules over the
sheaf of rings $R$. If $R$ is static, $D(R)$ is the $\infty$-categorical
enhancement of the usual unbounded derived category of the abelian category of
condensed $R$-modules (again in the naive sense: a sheaf of modules over a sheaf
of rings). For general $R$ one still has a $t$-structure on $D(R)$, but its heart
is no longer the module category once $R$ is not static. We use homological
grading conventions throughout.
\end{definition}

We can now give the definition. Below, $R^{\triangleright}$ (Clausen's notation:
the ``triangle'' denotes the underlying condensed ring) is a condensed animated
ring.

\begin{definition}[Analytic ring]
\label{def:1:analytic-ring}
An \emph{analytic ring} is a pair $R=(\ur{R},\,D(R))$ consisting of a condensed
(animated) ring $\ur{R}$ and a full subcategory $D(R)\subseteq D(\ur{R})$ --
its ``complete'' objects, the derived category of the analytic ring -- such
that:
\begin{enumerate}
\item $D(R)$ is closed under all limits and colimits (formed in $D(\ur{R})$);
\item if $N\in D(R)$ and $M\in D(\ur{R})$, then the internal hom
$\underline{R\mathrm{Hom}}_{\ur{R}}(M,N)$ lies in $D(R)$;
\item writing $\cpl{(-)}_R$ for the left adjoint to the inclusion
$D(R)\hookrightarrow D(\ur{R})$ (a completion functor), $\cpl{(-)}_R$ preserves
connectivity, i.e.\ sends $D(\ur{R})_{\geq 0}$ into $D(\ur{R})_{\geq 0}$; and
\item $\ur{R}\in D(R)$ (the unit, the underlying ring itself, is complete).
\end{enumerate}
\end{definition}

\begin{remark}[Existence of the completion functor]
\label{rmk:1:left-adjoint-automatic}
The left adjoint $\cpl{(-)}_R$ in axiom (3) exists automatically: given (1), the
inclusion of the presentable subcategory $D(R)$ admits a left adjoint with no
further set-theoretic hypothesis needed. Axiom (3) imposes that our analytic rings
are, in a suitable sense, connective, mirroring the connectivity of the ring.
The role of connectivity (axiom (3)) is to permit a reduction from a general
animated ring to its static truncation $\pi_0$.
\end{remark}

\begin{remark}[$A^{+}$, revisited]
\label{rmk:1:analytic-analog-of-Aplus}
The analytic ring structure -- the choice of $D(R)$ on top of $\ur{R}$ -- is
the analogue, in this theory, of the choice of a subring $A^{+}\subseteq A$ in
Huber's theory (\Cref{def:1:analytic-ring} versus the discussion of adic spaces
above).
\end{remark}

\begin{definition}[Map of analytic rings]
\label{def:1:map-analytic-rings}
A \emph{map of analytic rings} $R\to S$ is a map of condensed rings
$\ur{R}\to\ur{S}$ such that restriction of scalars along $\ur{R}\to\ur{S}$
carries $D(S)$ into $D(R)$: if $M\in D(S)$, then $M$ regarded as an
$\ur{R}$-module lies in $D(R)$.
\end{definition}

\begin{remark}[The $t$-structure on $D(R)$]
\label{rmk:1:t-structure}
There is always a $t$-structure on $D(R)$, described naively by intersection with
that of the enveloping category:
\begin{equation*}
D(R)_{\geq 0}=D(R)\cap D(\ur{R})_{\geq 0},\qquad
D(R)_{\leq 0}=D(R)\cap D(\ur{R})_{\leq 0},
\end{equation*}
so everything can be checked in the potentially more familiar $D(\ur{R})$. Its
heart
\[
D(R)^{\heartsuit}=D(R)\cap D(\ur{R})^{\heartsuit}
\]
is an abelian category. The modules are allowed to be $\mathbb{Z}$-graded, while
the ring is connective (positively graded). This abelian category already
determines the analytic ring: there is an equivalent axiomatics purely at the
abelian level, specifying $\ur{R}$ together with an abelian subcategory of
$D(\ur{R})^{\heartsuit}$ satisfying suitable axioms (one must, however, retain
axiom (4), that the unit is complete).
\end{remark}

A third, concrete perspective is in terms of free modules, i.e.\ spaces of
measures.

\begin{definition}[Completed free modules]
\label{def:1:free-modules}
For a profinite set $S$, the completed free module is
\begin{equation*}
R[S]:=\cpl{\bigl(\ur{R}[S]\bigr)}_R,
\end{equation*}
the completion of the free $\ur{R}$-module $\ur{R}[S]$ on $S$. These objects
generate $D(R)_{\geq 0}$ under colimits, and are the basic building blocks. There
is a further equivalent axiomatics taking as second datum not all of $D(R)$ but
just the collection of these free modules $R[S]\in D(\ur{R})$.
\end{definition}

\begin{remark}[Measures, and the completeness/integration interpretation]
\label{rmk:1:measures}
Intuitively $\ur{R}[S]$ is the space of \emph{finite} $\ur{R}$-linear
combinations of points of $S$ -- of Dirac measures on $S$ -- while its
completion $R[S]$ is a larger space of measures. In general $R[S]$ need not lie
in the heart (though in essentially all examples it does). An analytic ring
structure may thus be thought of as a choice of a space of measures on each
profinite set. Their role: for $M$ in the heart (say),
\begin{equation}
\label{eq:1:completeness-criterion}
M\in D(R)^{\heartsuit}\iff
\parbox{0.55\textwidth}{every map $f\colon\ur{R}[S]\to M$ of $\ur{R}$-modules
extends uniquely along $\ur{R}[S]\to R[S]$ to a map $R[S]\to M$.}
\end{equation}
Concretely, if $f$ is a function from $S$ to $M$ (a linear-algebra object with a
topology) and $\mu\in R[S]$ is one of the allowed measures, one obtains a
well-defined pairing $\int_S f\,d\mu\in M$. This is the sense in which $M$ is
``complete'': complete enough to integrate functions against the prescribed
class of measures, which is precisely the data of the analytic ring. (The
equivalence \eqref{eq:1:completeness-criterion} is by unwinding the left-adjoint
property.)
\end{remark}

\begin{remark}[Some monoidal structure not spelled out]
\label{rmk:1:monoidal}
$D(R)$ is symmetric monoidal, the completion functor $\cpl{(-)}_R$ is symmetric
monoidal, and the base-change (pullback) left adjoints associated to maps of
analytic rings are symmetric monoidal. The base-change functor -- the left
adjoint to the restriction-of-scalars in \Cref{def:1:map-analytic-rings} -- is
the most important functor of the theory; it is a completed tensor product,
obtained by tensoring up and then completing (tensoring up on the level of the
underlying condensed rings need not preserve $D(-)$). One of the theorems of the
theory is that the expected linear-algebra operations (symmetric powers, etc.)
carry over automatically; this is not obvious.
\end{remark}

\subsection{Colimits in analytic rings}

\begin{proposition}[Filtered and sifted colimits]
\label{prop:1:filtered-colimits}
Filtered colimits (more generally, sifted colimits) of analytic rings are
computed naively. For a filtered system $(R_i)_{i\in I}$,
\begin{equation*}
\Bigl(\varinjlim_i R_i\Bigr)^{\!\triangleright}=\varinjlim_{i\in I}\ur{R_i},
\qquad
\Bigl(\varinjlim_i R_i\Bigr)[S]=\varinjlim_i R_i[S]
\end{equation*}
on underlying rings and on free modules.
\end{proposition}

Pushouts are more interesting; they are the crucial construction, since they are
meant to produce the completed tensor products computing geometrically good fiber
products. Given maps of analytic rings from a base $K$ to $A$ and $B$, form the
pushout, written as a relative tensor product $A\otimes_K B$:
% board 34 @ 01:27:18 — /home/deven/documents/coding/vms/gpu/home/notetaker/output/dustin-clausen-124-analytic-stacks/boards/board-34.jpg
\[
\begin{tikzcd}
K \arrow[r] \arrow[d] & A \arrow[d] \\
B \arrow[r] & A\otimes_K B.
\end{tikzcd}
\]

\begin{proposition}[Derived category of a pushout]
\label{prop:1:pushout-derived-category}
Abstractly, $D(A\otimes_K B)$ is the full subcategory of
$D\bigl(\ur{A}\otimes_{\ur{K}}\ur{B}\bigr)$ consisting of those objects whose
underlying $\ur{A}$-module lies in $D(A)$ and whose underlying $\ur{B}$-module
lies in $D(B)$. (Equivalently, it is the relative tensor product
$D(A)\otimes_{D(K)} D(B)$ of module categories.)
\end{proposition}

\begin{remark}[The underlying ring must be completed]
\label{rmk:1:pushout-completion}
Caution: the pair
\begin{equation}
\label{eq:1:pushout-not-analytic}
\bigl(\ur{A}\otimes_{\ur{K}}\ur{B},\ D(A\otimes_K B)\bigr)
\end{equation}
satisfies axioms (1)--(3) of \Cref{def:1:analytic-ring} but not axiom (4): it is
almost an analytic ring, the only failure being that the unit
$\ur{A}\otimes_{\ur{K}}\ur{B}$ is not complete. One repairs this by applying the
completion, which changes the underlying ring but leaves the category $D(A\otimes_K B)$
unchanged; that the result is again a ring must be proved and is not obvious. In
full generality the completion can be hard to understand: abstractly one iterates
the completions for $A$ and for $B$, sandwiched, countably many times, and takes a
colimit. In practice, however, it can be computed, and it produces the
geometrically correct completed tensor products of analytic geometry.
\end{remark}

\begin{remark}[Trivial (maximal) analytic ring structure]
\label{rmk:1:trivial-structure}
Every condensed ring $\ur{A}$ carries the trivial analytic ring structure
$D(A):=D(\ur{A})$ (everything is complete); write this analytic ring simply as
$\ur{A}$. With this convention the notation in \eqref{eq:1:pushout-not-analytic}
is consistent, and every condensed ring is an analytic ring with its maximal
structure.
\end{remark}

\subsection{Example: solid analytic rings}

We turn to an example, connected to adic spaces (Huber pairs): the solid analytic
rings, which encode a non-Archimedean condition. Given an analytic ring $R$, it
is illuminating to look at the free modules on profinite sets to see which spaces
of measures occur. In the linear setting, the natural object built from the basic
profinite set $\mathbb{N}\cup\{\infty\}$ is the following.

\begin{definition}[Measures on $\mathbb{N}$]
\label{def:1:measures-on-N}
For an analytic ring $R$, set
\begin{equation*}
\Mmeas_R(\mathbb{N}):=R[\mathbb{N}\cup\{\infty\}]\big/\,R[\{\infty\}],
\end{equation*}
the completed free module on the one-point compactification
$\mathbb{N}\cup\{\infty\}$, modulo the free module on the point $\infty$ (which is
just $R$). It is the space of measures on $\mathbb{N}$; dually it classifies null
sequences in $R$-modules.
\end{definition}

\begin{remark}[The ring structure and the variable $T$]
\label{rmk:1:MRN-ring}
Addition on $\mathbb{N}$ induces a ring structure on $\Mmeas_R(\mathbb{N})$; it is
not hard to show this. Writing the resulting element as $T$, the ring
$\Mmeas_R(\mathbb{N})$ sits between the two extremes
\begin{equation}
\label{eq:1:MRN-between}
R[T]\longrightarrow \Mmeas_R(\mathbb{N})\longrightarrow R[[T]],
\end{equation}
the polynomial algebra and the power series algebra over $R$ -- as one expects
of a space of null sequences (really its dual, some summability condition).
Geometrically one has the affine line, a version of the formal neighbourhood of
the origin, and $\Mmeas_R(\mathbb{N})$ somewhere in between. (In the universal
case $R=\Zsolid$ this ring lives in degree $0$; in general it is a summand and
need not lie in the heart, but by base change from the universal case it is a
ring in whatever context one wants.)
\end{remark}

\begin{definition}[Solid analytic ring]
\label{def:1:solid}
An analytic ring $R$ is \emph{solid} if
\begin{equation*}
\Mmeas_R(\mathbb{N})\big/(T-1)=0,
\end{equation*}
i.e.\ if $\Mmeas_R(\mathbb{N})$, viewed as sitting inside the affine line via
$T$, avoids the point $1$; equivalently, multiplication by $T-1$ on
$\Mmeas_R(\mathbb{N})$ is an isomorphism. Morally, this places
$\Mmeas_R(\mathbb{N})$ inside the open unit disk, a non-Archimedean condition.
\end{definition}

\begin{remark}[Solidity as summability of null sequences]
\label{rmk:1:solid-summable}
Solidity is equivalent to the existence of a (necessarily unique) measure
\begin{equation*}
\mu\in\Mmeas_R(\mathbb{N})\quad\text{with}\quad (T-1)\mu=1,
\end{equation*}
where $1\in\Mmeas_R(\mathbb{N})$ is the unit, i.e.\ the sequence
$(1,0,0,\dots)$. Under \eqref{eq:1:MRN-between} the measure $\mu$ maps to
$\sum_{n\geq 0} T^{n}$ in the power series ring. Since $\Mmeas_R(\mathbb{N})$
classifies null sequences, and $\mu$ pairs with a null sequence by taking it as
the coefficients of $\sum_n c_n T^n$ and setting $T=1$, the interpretation is
that \emph{every null sequence is summable} -- the classical non-Archimedean
condition.
\end{remark}

Clausen then stated (without full detail) the fundamental example.

\begin{theorem}[The solid analytic ring $\Zsolid$]
\label{thm:1:Zsolid}
There exists a solid analytic ring
\[
\Zsolid=\bigl(\mathbb{Z},\,D(\Zsolid)\bigr),
\]
whose underlying condensed ring is the usual integers $\mathbb{Z}$ (with discrete
topology), such that an analytic ring $R$ is solid if and only if there exists a
(necessarily unique) map of analytic rings $\Zsolid\to R$.
\end{theorem}

The category $D(\Zsolid)$ can be understood very explicitly; the following are
the basic facts of linear algebra in it.

\begin{proposition}[Structure of $D(\Zsolid)$]
\label{prop:1:Zsolid-structure}
Let $S=\varprojlim_n S_n$ be a profinite set.
\begin{enumerate}
\item The free module is
\begin{equation}
\label{eq:1:Zsolid-free-module}
\Zsolid[S]=\varprojlim_n \mathbb{Z}[S_n]=\varprojlim_n \mathbb{Z}^{S_n},
\end{equation}
a countable inverse limit of finite free modules, abstractly isomorphic to a
(countably infinite) product $\prod_I\mathbb{Z}$ unless $S$ is finite.
\item These products $\prod_I\mathbb{Z}$ are compact projective generators of the
heart $D(\Zsolid)^{\heartsuit}$ and lie in degree $0$; consequently $D(\Zsolid)$
is the usual derived category of its heart. They are flat for the completed
tensor product $-\otimes_{\Zsolid}-$, and tensor products are computed by the
infinite distributive law
\begin{equation*}
\Bigl(\prod_I\mathbb{Z}\Bigr)\otimes_{\Zsolid}\Bigl(\prod_J\mathbb{Z}\Bigr)
=\prod_{I\times J}\mathbb{Z}.
\end{equation*}
\item The finitely presented objects of $D(\Zsolid)^{\heartsuit}$ (which generate
it under filtered colimits) form an abelian subcategory closed under extensions,
and every finitely presented $M$ admits a resolution by products
$\prod_I\mathbb{Z}$ of length $\leq 2$ (a complex with three nonzero terms and two
nonzero maps).
\end{enumerate}
\end{proposition}

\begin{remark}[Lightness is used]
\label{rmk:1:lightness-used}
The flatness in \Cref{prop:1:Zsolid-structure}(2) is a place where the light
restriction matters: Efimov proved that it fails if one increases the
cardinalities of the profinite sets.
\end{remark}

\begin{remark}[$\Zsolid$ as a regular ring of dimension $2$]
\label{rmk:1:Zsolid-regular-dim-2}
The length-$\leq 2$ resolutions of \Cref{prop:1:Zsolid-structure}(3) say that
$\Zsolid$ behaves like a regular ring of dimension $2$: whereas $\mathbb{Z}$ is
regular of dimension $1$, one picks up an extra dimension, attributable to the
non-Hausdorff phenomena occurring in solid abelian groups. All told, one gets a
very good handle on this category.
\end{remark}

\begin{remark}[A single generator]
\label{rmk:1:single-generator}
In \eqref{eq:1:Zsolid-free-module} the index set $I$ is countable, and in fact a
single generator suffices: the Cantor set $\{0,1\}^{\mathbb{N}}$ always generates,
a general phenomenon.
\end{remark}

\subsection{Closing remarks}

Clausen closed with two points raised in discussion. First, an application: the
theory yields very general Riemann--Roch theorems in analytic geometry, using
these derived categories and the flexibility of the formalism in an essential
way. Second, on the choice of framework: one uses light
condensed sets in preference to general condensed sets to avoid choosing a
cardinality bound; with a strong-limit cardinal one does get compactly generated
derived categories, which is psychologically comforting though not so important in
practice. Both $D(\Zsolid)$ and its heart are $\aleph_1$-compactly generated, though not
compactly generated -- a consequence of the light restriction; Clausen thought,
but was not certain, that $D(\Zsolid)$ is not dualizable. Finally, the reason for choosing $\mathbb{E}_\infty$ or animated
commutative rings, rather than $\mathbb{E}_1$ or $\mathbb{E}_2$, is that for them
coproducts coincide with relative tensor products; moving to settings where this
fails is a real inconvenience.
